\documentclass{resume}

\usepackage[left=0.75in,top=0.50in,right=0.75in,bottom=0.50in]{geometry}
\newcommand{\tab}[1]{\hspace{.1667\textwidth}\rlap{#1}}
\newcommand{\itab}[1]{\hspace{0em}\rlap{#1}}
\name{Kendrake Tsui}
\address{(858) 610-4089 • tsui.kendrake@gmail.com • linkedin.com/in/kendraketsui • github.com/kaytsui}

\begin{document}

\begin{rSection}{Education}

\begin{rSubsection}{University of California, Los Angeles}{Expected June 2022}{Bachelor of Science in Computer Science}{GPA: 3.91}
\item \textit{\textbf{Relevant Coursework:} Data Structures \& Algorithms, Operating Systems, Programming Languages, Machine Learning, Artificial Intelligence, Databases, Web Development, Computer Networking (IP)}

\end{rSubsection}

\end{rSection}


\begin{rSection}{Work Experience}

\begin{rSubsection}{San Diego Zoo Wildlife Alliance}{San Diego, CA}{Software Engineering Intern}{June 2021 - September 2021}
\item Created an interactive dashboard using Python's Dash framework to dynamically render images taken by camera traps, allowing for in time remote animal tracking to aid population sustainability efforts.
\item Conducted meetings with end users to determine core features of the dashboard that include a data \\ visualization page, an image history page, and an object recognition labeling component.
\end{rSubsection}

\begin{rSubsection}{Code.org}{Remote Work}{Software Engineering and Curriculum  Intern}{February 2021 - June 2021}
\item Wrote validation code in JavaScript for 35+ levels of Code.org's CS Fundamentals curriculum to give more structured feedback on level completion for 750,000+ upper elementary students nationwide.
\item Refactored code to streamline validation structure and collaborated with a curriculum writer and engineer to create a modular validation template to be used by future curriculum interns.
\end{rSubsection}
\end{rSection}

\begin{rSection}{Projects}
\begin{rSubsection}{Harmony Project Mobile App for Music Education}{Los Angeles, CA}{Full Stack Developer}{November 2020 - February 2021}
\item Developed a React Native mobile app in a team of 8 for Harmony Project, a 501(c)(3) nonprofit focused on providing equitable music education for K-12 students in Los Angeles.
\item Implemented the front-end and back-end of the student assignment page and file attachment interface.
\end{rSubsection}

\begin{rSubsection}{Project Horizon}{Los Angeles, CA}{Founder and Developer}{July 2020 - December 2020}
\item Managed, innovated, and developed a website to elevate faculty perspectives, create a diverse and inclusive multi-disciplinary forum, and broaden student awareness of subjects outside their major.
\item Led a cross-functional team of three developers and UX designers to select website color palette, theme, logo, and design to symbolize Plato’s Allegory of the cave and UCLA’s Fiat Lux motto.
\end{rSubsection}
\end{rSection}

\begin{rSection}{Leadership}
\begin{rSubsection}{Upsilon Pi Epsilon (Top 1/3 of Computer Science Class)}{Los Angeles, CA}{Advocacy Director}{June 2020 - Present}
\item Design an inaugural Zoom welcome event in collaboration with 6 engineering organizations to support new computer science students in their transition to UCLA, attended by over 235 new students.
\item Spearhead a collaboration between the Computer Science Faculty and Education Department Faculty to discuss best practices in how faculty can integrate diversity, inclusion, and ethics into their curriculum.
\end{rSubsection}
\end{rSection}



\begin{rSection}{Technical Skills}

\begin{tabular}{ @{} >{\bfseries}l @{\hspace{6ex}} l }
Programming Languages & C, C++, Python, JavaScript, Node.js, TypeScript, Java, Bash \\
Technology and Frameworks & Git, React, Dash, SQL, Firebase, Linux  \\
\end{tabular}

\end{rSection}

\end{document}
